%% Experimental Setup Section
\section{Experimental Setup}
\label{sec:experiments}

\subsection{QUASAR v14 Framework}

QUASAR v14 (Quantum Universal Adaptive State pReparation) is a multi-target
quantum state preparation framework built on the following components:

\paragraph{Base RL algorithm.}
We use Soft Actor-Critic (SAC) \cite{haarnoja2018soft} with automatic entropy
tuning. The actor and critic networks are multilayer perceptrons with two hidden
layers of 256 units each and ReLU activations. The replay buffer has capacity
$10^6$ transitions, and we use a batch size of 256 with a learning rate of
$3 \times 10^{-4}$ for all networks.

\paragraph{Continual multi-target learning.}
To enable a single agent to learn control policies for multiple entangled state
families without catastrophic forgetting, we augment SAC with Dark Experience
Replay++ (DER++) \cite{buzzega2020dark}. DER++ stores a buffer of (state, logit)
pairs from past training trajectories and adds a distillation loss term to the
SAC objective:
\begin{equation}
  \mathcal{L}_{\mathrm{DER++}} = \mathcal{L}_{\mathrm{SAC}} +
  \lambda \mathbb{E}_{(s, z) \sim \mathcal{B}_{\mathrm{DER}}}
  \bigl[\|Q(s, \cdot) - z\|^2\bigr],
\end{equation}
where $\mathcal{B}_{\mathrm{DER}}$ is the DER++ buffer and $\lambda = 0.5$.

\paragraph{Exploration enhancement.}
We introduce Stochastic Layer Modulation (SLM), a novel exploration mechanism
that randomly scales the activations of a randomly selected hidden layer by a
factor drawn from $\mathcal{U}(0.8, 1.2)$ during exploration episodes. SLM
provides structured perturbations to the policy that are more informative than
isotropic Gaussian noise, particularly near the target state where the gradient
landscape is flat.

\paragraph{Hyperparameter optimisation.}
The entropy coefficient $\alpha$ is optimised using DEHB \cite{awad2021dehb}
with 12 trials per qubit level. Each trial runs for an initial budget of 40,000
steps, with ACC providing early termination for unreachable configurations.
The best $\alpha$ from the HP search phase is used for the full-budget run,
with the budget set to $T^* \times 1.2$ as predicted by ACC.

\subsection{Quantum Simulation Environment}

We simulate $n$-qubit systems using SiliQun v2 \cite{TODO_siliqun}, a
silicon spin qubit simulator that models exchange-coupled spin-1/2 particles
with realistic device parameters. The Hamiltonian is:
\begin{equation}
  H_0 = \sum_{i=1}^{n} \frac{\omega_i}{2} \sigma_z^{(i)}
        + \sum_{\langle i,j\rangle} J_{ij} \bigl(\sigma_x^{(i)}\sigma_x^{(j)}
        + \sigma_y^{(i)}\sigma_y^{(j)}\bigr),
\end{equation}
where $\omega_i / 2\pi \in [4.5, 5.5]\,\mathrm{GHz}$ are qubit frequencies and
$J_{ij} / 2\pi \in [1, 10]\,\mathrm{MHz}$ are exchange coupling strengths,
drawn from the SiMOS device profile \cite{TODO_siliqun}.

\subsection{Baseline Methods}

We compare ACC against three baselines:

\begin{enumerate}
  \item \textbf{Fixed budget (FB)}: the standard approach used in prior work
    \cite{yao2021reinforcement,moro2021quantum}, with a fixed budget of
    $500{,}000 \times 1.5^{n-2}$ steps per qubit level.
  \item \textbf{Plateau counter (PC)}: the heuristic used in QUASAR v13,
    which stops training when the best fidelity has not improved for 20,000
    consecutive steps.
  \item \textbf{ASHA}: the Asynchronous Successive Halving Algorithm
    \cite{li2020system} applied to the training budget allocation problem,
    using the same 12-trial budget as DEHB+ACC.
\end{enumerate}

\subsection{Evaluation Protocol}

All experiments are run with three random seeds (42, 123, 456) per
(target family, qubit count) combination. We report the mean and standard
deviation of the best fidelity achieved, the total number of training steps,
and the wall-clock time on the Aziz Supercomputer A100 nodes. Statistical
significance is assessed using a two-sided Wilcoxon signed-rank test with
$\alpha = 0.05$.

\subsection{Computational Resources}

All experiments were performed on the Aziz Supercomputer at King Abdulaziz
University, using NVIDIA A100-PCIE-40GB GPUs (job 181423) and
NVIDIA H100-SXM5-80GB GPUs (job 181404). Each job ran under PBS Pro with
a 72-hour walltime limit. The total compute budget for the main experiments
is approximately 200 GPU-hours.
